
\documentclass{article}
\usepackage{colortbl}
\usepackage{makecell}
\usepackage{multirow}
\usepackage{supertabular}

\begin{document}

\newcounter{utterance}

\centering \large Interaction Transcript for game `dond', experiment `coop\_en', episode 13 with reflexion{-}Qwen3.5{-}4B.
\vspace{24pt}

{ \footnotesize  \setcounter{utterance}{1}
\setlength{\tabcolsep}{0pt}
\begin{supertabular}{c@{$\;$}|p{.15\linewidth}@{}p{.15\linewidth}p{.15\linewidth}p{.15\linewidth}p{.15\linewidth}p{.15\linewidth}}
    \# & $\;$A & \multicolumn{4}{c}{Game Master} & $\;\:$B\\
    \hline

    \theutterance \stepcounter{utterance}  
    & & \multicolumn{4}{p{0.6\linewidth}}{
        \cellcolor[rgb]{0.9,0.9,0.9}{
            \makecell[{{p{\linewidth}}}]{
                \texttt{\tiny{[P1$\langle$GM]}}
                \texttt{You have attempted this same task before and did not succeed. These are the plans you wrote after those attempts.} \\
\texttt{Trial \#1: New plan: I failed because I assumed the other player would accept my offer of combs for a statue and a stamp, but I did not verify if their secret valuation of the combs was high enough to make that trade mutually beneficial. Since the goal is to maximize the *sum* of both scores, I should have calculated the total value of all items first (statue: 1, combs: 8, stamp: 1, total: 10) and then proposed a split that divides this total as evenly as possible (e.g., 5 each), rather than focusing solely on my own high{-}value items. Next time, I will first calculate the total value of the set, then propose a division where the sum of my items plus the sum of the other player's items equals the total, ensuring the split is roughly 50/50 or the most equitable distribution possible, regardless of which specific items I hold.} \\
\\ 
\\ 
\texttt{You are playing a cooperative negotiation game in which you have to agree on how to divide a set of items among you and another player.} \\
\\ 
\texttt{Rules:} \\
\texttt{(a) You and the other player are given a set of items. Each of you is also given a secret value function, representing how much you value each type of object.} \\
\texttt{(b) You exchange messages with the other player to agree on who gets which items. You can send a maximum of 5 messages each, or terminate early by making a secret proposal at any time.} \\
\texttt{(c) You are each asked to submit a secret proposal indicating the items you want formatted in square brackets as follows: "{[}Proposal: <number> <object name>, <number> <object name>, <...>{]}"} \\
\texttt{(d) If your proposals are complementary, i.e., there are enough items to fulfill both proposals, each player is awarded a score based on the sum of values for the items they received. Otherwise, both of you get zero points.} \\
\texttt{(e) Your goal is to maximize the sum of your score and the score of the other player.} \\
\\ 
\texttt{Let us start.} \\
\\ 
\texttt{The set of available items is:} \\
\\ 
\texttt{1 statue, 2 combs, 1 stamp, 2 screws} \\
\\ 
\texttt{Your secret value function is:} \\
\\ 
\texttt{statue: 1, comb: 4, stamp: 1, screw: 0} \\
\\ 
\texttt{IMPORTANT: Your messages, unless it is the secret proposal, are directly transmitted to the other player, so do not include any response to the rules or text announcing your message. To make a secret proposal, use the indicated format. Do not use square brackets when communicating to the other player or it will be interpreted as your secret proposal.} \\
            }
        }
    }
    & \\ \\

    \theutterance \stepcounter{utterance}  
    & \multicolumn{4}{p{0.6\linewidth}}{
        \cellcolor[rgb]{0.9,0.9,0.9}{
            \makecell[{{p{\linewidth}}}]{
                \texttt{\tiny{[P1$\rangle$GM]}}
                \texttt{I will propose a split that maximizes the total value. My total value is 1 (statue) + 2*4 (combs) + 1 (stamp) + 2*0 (screws) = 10. The other player likely values the combs highly as well (since they are the most numerous and usually high{-}value in these puzzles, or at least the statue and stamp). To ensure a successful trade where both get points, I should offer the items I value least (statue and stamp) to the other player in exchange for the items they value most (combs), while keeping the combs for myself.} \\
\\ 
\texttt{Let's assume the other player values combs highly. If I give them the statue and stamp (value 2 to me, likely high to them) and they give me the combs (value 8 to me, likely high to them), we both get a good score.} \\
\texttt{My plan: I keep the 2 combs (my value 8). I give the statue and stamp to the other player. I take the screws (value 0) as a "cost" to make the deal work, or perhaps the other player takes the screws.} \\
\texttt{Actually, let's look at the items:} \\
\texttt{Me: Statue(1), Comb(4), Stamp(1), Screw(0). Total = 10.} \\
\texttt{Other: Statue(1), Comb(4), Stamp(1), Screw(0). Total = 10. (Assuming symmetric values for simplicity, or at least they value combs).} \\
\\ 
\texttt{If I propose:} \\
\texttt{Me: 2 combs, 2 screws. (My value: 8 + 0 = 8).} \\
\texttt{Other: 1 statue, 1 stamp. (Their value: 1 + 1 = 2).} \\
\texttt{Total sum = 10. This is bad for the other player.} \\
\\ 
\texttt{Better split:} \\
\texttt{Me: 2 combs, 1 statue. (My value: 8 + 1 = 9).} \\
\texttt{Other: 1 stamp, 2 screws. (Their value: 1 + 0 = 1).} \\
\texttt{Total sum = 10. Still bad.} \\
\\ 
\texttt{Wait, the goal is to maximize the *sum*. The sum is fixed at 20 if we split everything perfectly. The constraint is that we must agree on a split. If I propose a split where the other player gets nothing they value, they might reject it (or if the game forces a secret proposal, it might fail if not complementary).} \\
\texttt{Let's try to give the other player the items I value least, which are the screws and maybe the statue/stamp.} \\
\texttt{If I give them the 2 screws and 1 stamp. Their value is likely low.} \\
\texttt{If I give them the 1 statue and 2 screws. Their value is likely low.} \\
\\ 
\texttt{Let's} \\
            }
        }
    }
    & & \\ \\

    \theutterance \stepcounter{utterance}  
    & & & \multicolumn{2}{p{0.3\linewidth}}{
        \cellcolor[rgb]{0.9,0.9,0.9}{
            \makecell[{{p{\linewidth}}}]{
                \texttt{\tiny{[GM$|$GM]}}
                \texttt{continue} \\
            }
        }
    }
    & & \\ \\

    \theutterance \stepcounter{utterance}  
    & & & \multicolumn{2}{p{0.3\linewidth}}{
        \cellcolor[rgb]{0.9,0.9,0.9}{
            \makecell[{{p{\linewidth}}}]{
                \texttt{\tiny{[GM$|$GM]}}
                \texttt{I will propose a split that maximizes the total value. My total value is 1 (statue) + 2*4 (combs) + 1 (stamp) + 2*0 (screws) = 10. The other player likely values the combs highly as well (since they are the most numerous and usually high{-}value in these puzzles, or at least the statue and stamp). To ensure a successful trade where both get points, I should offer the items I value least (statue and stamp) to the other player in exchange for the items they value most (combs), while keeping the combs for myself.} \\
\\ 
\texttt{Let's assume the other player values combs highly. If I give them the statue and stamp (value 2 to me, likely high to them) and they give me the combs (value 8 to me, likely high to them), we both get a good score.} \\
\texttt{My plan: I keep the 2 combs (my value 8). I give the statue and stamp to the other player. I take the screws (value 0) as a "cost" to make the deal work, or perhaps the other player takes the screws.} \\
\texttt{Actually, let's look at the items:} \\
\texttt{Me: Statue(1), Comb(4), Stamp(1), Screw(0). Total = 10.} \\
\texttt{Other: Statue(1), Comb(4), Stamp(1), Screw(0). Total = 10. (Assuming symmetric values for simplicity, or at least they value combs).} \\
\\ 
\texttt{If I propose:} \\
\texttt{Me: 2 combs, 2 screws. (My value: 8 + 0 = 8).} \\
\texttt{Other: 1 statue, 1 stamp. (Their value: 1 + 1 = 2).} \\
\texttt{Total sum = 10. This is bad for the other player.} \\
\\ 
\texttt{Better split:} \\
\texttt{Me: 2 combs, 1 statue. (My value: 8 + 1 = 9).} \\
\texttt{Other: 1 stamp, 2 screws. (Their value: 1 + 0 = 1).} \\
\texttt{Total sum = 10. Still bad.} \\
\\ 
\texttt{Wait, the goal is to maximize the *sum*. The sum is fixed at 20 if we split everything perfectly. The constraint is that we must agree on a split. If I propose a split where the other player gets nothing they value, they might reject it (or if the game forces a secret proposal, it might fail if not complementary).} \\
\texttt{Let's try to give the other player the items I value least, which are the screws and maybe the statue/stamp.} \\
\texttt{If I give them the 2 screws and 1 stamp. Their value is likely low.} \\
\texttt{If I give them the 1 statue and 2 screws. Their value is likely low.} \\
\\ 
\texttt{Let's} \\
            }
        }
    }
    & & \\ \\

    \theutterance \stepcounter{utterance}  
    & & \multicolumn{4}{p{0.6\linewidth}}{
        \cellcolor[rgb]{0.9,0.9,0.9}{
            \makecell[{{p{\linewidth}}}]{
                \texttt{\tiny{[GM$\rangle$P2]}}
                \texttt{You have attempted this same task before and did not succeed. These are the plans you wrote after those attempts.} \\
\texttt{Trial \#1: New plan: I failed because I submitted a proposal containing all items, which is not complementary to any possible proposal from the other player (since they must also receive items to get a non{-}zero score). To succeed, I must propose a subset of items that leaves enough items for the other player to form a valid, complementary set. Given my values (statue: 4, stamp: 2, combs: 1, screws: 1), I should prioritize the statue and stamp. I will propose taking the statue and one comb, leaving the stamp, remaining comb, and two screws for the other player. This ensures both players have items, maximizing the chance of a successful agreement while securing my highest{-}value items.} \\
\\ 
\\ 
\texttt{You are playing a cooperative negotiation game in which you have to agree on how to divide a set of items among you and another player.} \\
\\ 
\texttt{Rules:} \\
\texttt{(a) You and the other player are given a set of items. Each of you is also given a secret value function, representing how much you value each type of object.} \\
\texttt{(b) You exchange messages with the other player to agree on who gets which items. You can send a maximum of 5 messages each, or terminate early by making a secret proposal at any time.} \\
\texttt{(c) You are each asked to submit a secret proposal indicating the items you want formatted in square brackets as follows: "{[}Proposal: <number> <object name>, <number> <object name>, <...>{]}"} \\
\texttt{(d) If your proposals are complementary, i.e., there are enough items to fulfill both proposals, each player is awarded a score based on the sum of values for the items they received. Otherwise, both of you get zero points.} \\
\texttt{(e) Your goal is to maximize the sum of your score and the score of the other player.} \\
\\ 
\texttt{Let us start.} \\
\\ 
\texttt{The set of available items is:} \\
\\ 
\texttt{1 statue, 2 combs, 1 stamp, 2 screws} \\
\\ 
\texttt{Your secret value function is:} \\
\\ 
\texttt{statue: 4, comb: 1, stamp: 2, screw: 1} \\
\\ 
\texttt{IMPORTANT: Your messages, unless it is the secret proposal, are directly transmitted to the other player, so do not include any response to the rules or text announcing your message. To make a secret proposal, use the indicated format. Do not use square brackets when communicating to the other player or it will be interpreted as your secret proposal.} \\
\\ 
\\ 
\texttt{I will propose a split that maximizes the total value. My total value is 1 (statue) + 2*4 (combs) + 1 (stamp) + 2*0 (screws) = 10. The other player likely values the combs highly as well (since they are the most numerous and usually high{-}value in these puzzles, or at least the statue and stamp). To ensure a successful trade where both get points, I should offer the items I value least (statue and stamp) to the other player in exchange for the items they value most (combs), while keeping the combs for myself.} \\
\\ 
\texttt{Let's assume the other player values combs highly. If I give them the statue and stamp (value 2 to me, likely high to them) and they give me the combs (value 8 to me, likely high to them), we both get a good score.} \\
\texttt{My plan: I keep the 2 combs (my value 8). I give the statue and stamp to the other player. I take the screws (value 0) as a "cost" to make the deal work, or perhaps the other player takes the screws.} \\
\texttt{Actually, let's look at the items:} \\
\texttt{Me: Statue(1), Comb(4), Stamp(1), Screw(0). Total = 10.} \\
\texttt{Other: Statue(1), Comb(4), Stamp(1), Screw(0). Total = 10. (Assuming symmetric values for simplicity, or at least they value combs).} \\
\\ 
\texttt{If I propose:} \\
\texttt{Me: 2 combs, 2 screws. (My value: 8 + 0 = 8).} \\
\texttt{Other: 1 statue, 1 stamp. (Their value: 1 + 1 = 2).} \\
\texttt{Total sum = 10. This is bad for the other player.} \\
\\ 
\texttt{Better split:} \\
\texttt{Me: 2 combs, 1 statue. (My value: 8 + 1 = 9).} \\
\texttt{Other: 1 stamp, 2 screws. (Their value: 1 + 0 = 1).} \\
\texttt{Total sum = 10. Still bad.} \\
\\ 
\texttt{Wait, the goal is to maximize the *sum*. The sum is fixed at 20 if we split everything perfectly. The constraint is that we must agree on a split. If I propose a split where the other player gets nothing they value, they might reject it (or if the game forces a secret proposal, it might fail if not complementary).} \\
\texttt{Let's try to give the other player the items I value least, which are the screws and maybe the statue/stamp.} \\
\texttt{If I give them the 2 screws and 1 stamp. Their value is likely low.} \\
\texttt{If I give them the 1 statue and 2 screws. Their value is likely low.} \\
\\ 
\texttt{Let's} \\
            }
        }
    }
    & \\ \\

    \theutterance \stepcounter{utterance}  
    & & & \multicolumn{4}{p{0.6\linewidth}}{
        \cellcolor[rgb]{0.9,0.9,0.9}{
            \makecell[{{p{\linewidth}}}]{
                \texttt{\tiny{[GM$\langle$P2]}}
                \texttt{{[}Proposal: 1 statue, 1 stamp, 2 screws{]}} \\
            }
        }
    }
     \\ \\

    \theutterance \stepcounter{utterance}  
    & & & \multicolumn{2}{p{0.3\linewidth}}{
        \cellcolor[rgb]{0.9,0.9,0.9}{
            \makecell[{{p{\linewidth}}}]{
                \texttt{\tiny{[GM$|$GM]}}
                \texttt{proposal} \\
            }
        }
    }
    & & \\ \\

    \theutterance \stepcounter{utterance}  
    & & & \multicolumn{2}{p{0.3\linewidth}}{
        \cellcolor[rgb]{0.9,0.9,0.9}{
            \makecell[{{p{\linewidth}}}]{
                \texttt{\tiny{[GM$|$GM]}}
                [1, 0, 1, 2]
            }
        }
    }
    & & \\ \\

    \theutterance \stepcounter{utterance}  
    & & \multicolumn{4}{p{0.6\linewidth}}{
        \cellcolor[rgb]{0.9,0.9,0.9}{
            \makecell[{{p{\linewidth}}}]{
                \texttt{\tiny{[P1$\langle$GM]}}
                \texttt{The other player has submitted a secret proposal. It is now your turn to submit a proposal of your own.} \\
            }
        }
    }
    & \\ \\

    \theutterance \stepcounter{utterance}  
    & \multicolumn{4}{p{0.6\linewidth}}{
        \cellcolor[rgb]{0.9,0.9,0.9}{
            \makecell[{{p{\linewidth}}}]{
                \texttt{\tiny{[P1$\rangle$GM]}}
                \texttt{{[}Proposal: 2 comb, 1 statue{]}} \\
            }
        }
    }
    & & \\ \\

    \theutterance \stepcounter{utterance}  
    & & & \multicolumn{2}{p{0.3\linewidth}}{
        \cellcolor[rgb]{0.9,0.9,0.9}{
            \makecell[{{p{\linewidth}}}]{
                \texttt{\tiny{[GM$|$GM]}}
                \texttt{proposal} \\
            }
        }
    }
    & & \\ \\

    \theutterance \stepcounter{utterance}  
    & & & \multicolumn{2}{p{0.3\linewidth}}{
        \cellcolor[rgb]{0.9,0.9,0.9}{
            \makecell[{{p{\linewidth}}}]{
                \texttt{\tiny{[GM$|$GM]}}
                [1, 2, 0, 0]
            }
        }
    }
    & & \\ \\

    \theutterance \stepcounter{utterance}  
    & & & \multicolumn{2}{p{0.3\linewidth}}{
        \cellcolor[rgb]{0.9,0.9,0.9}{
            \makecell[{{p{\linewidth}}}]{
                \texttt{\tiny{[GM$|$GM]}}
                [[1, 2, 0, 0], [1, 0, 1, 2]]
            }
        }
    }
    & & \\ \\

    \theutterance \stepcounter{utterance}  
    & \multicolumn{4}{p{0.6\linewidth}}{
        \cellcolor[rgb]{0.95,0.95,0.95}{
            \makecell[{{p{\linewidth}}}]{
                \texttt{\tiny{[P2$\rightarrow$P1]}}
                \texttt{Post{-}episode Reflexion (not sent to GM):} \\
\texttt{New plan: I failed because I proposed a split (2 combs, 1 statue) that left the other player with only 1 stamp and 2 screws. Even if the other player values combs highly, they received items I valued at 0 (screws) and low value (stamp), likely resulting in a total score for them that was too low to make the deal mutually beneficial or complementary. The goal is to maximize the *sum* of both scores, which means finding a split where the other player receives enough value to accept. I should have calculated the total value of all items (2 statues, 4 combs, 1 stamp, 2 screws = 20 total value assuming symmetric values, or at least ensuring the other player gets a significant portion). Next time, I will propose a split where the other player receives the items I value least (screws) plus one or two other items (like the stamp or one comb) to ensure their total value is high enough to make the agreement successful, rather than giving them only the lowest{-}value items. Specifically, I will propose: {[}Proposal: 2 comb, 1 screw{]} for myself, leaving the other player with 1 statue, 1 comb, 1 stamp, and 1 screw (or similar balanced distribution) to ensure the sum is maximized and the deal is accepted.} \\
            }
        }
    }
    & & \\ \\

    \theutterance \stepcounter{utterance}  
    & & & \multicolumn{4}{p{0.6\linewidth}}{
        \cellcolor[rgb]{0.95,0.95,0.95}{
            \makecell[{{p{\linewidth}}}]{
                \texttt{\tiny{[P1$\rightarrow$P2]}}
                \texttt{Post{-}episode Reflexion (not sent to GM):} \\
\texttt{New plan: I failed because I proposed taking all the items I valued (statue, stamp, and screws), leaving the other player with nothing. Since the other player also needs items to receive a non{-}zero score, my proposal was not complementary to any valid response from them. To succeed, I must propose a split where I take a subset of items and leave a complementary subset for the other player. Given my values (statue: 4, stamp: 2, combs: 1, screws: 1), I should prioritize the statue and stamp. I will propose taking the statue and one comb (my value: 5), leaving the stamp, one comb, and two screws for the other player. This ensures both players have items, maximizing the chance of a successful agreement while securing my highest{-}value items.} \\
            }
        }
    }
     \\ \\

\end{supertabular}
}

\end{document}
